% Recovering Time-Varying Correlation with DCC-GARCH.
% Controlled study of correlation tracking and portfolio-risk misstatement.
% Every number in this manuscript is verified against results/results.json by
% scripts/check_paper_numbers.py; do not edit numbers without re-running it.
\documentclass[11pt]{article}
\usepackage[margin=1in]{geometry}
\usepackage{amsmath,amssymb}
\usepackage{booktabs}
\usepackage[round]{natbib}
\usepackage{microtype}
\usepackage[hidelinks]{hyperref}
\setlength{\emergencystretch}{2em}

\newcommand{\E}{\mathbb{E}}
\newcommand{\Var}{\mathrm{Var}}
\newcommand{\Cov}{\mathrm{Cov}}
\newcommand{\diag}{\mathrm{diag}}
\newcommand{\Rt}{R_{t}}
\newcommand{\Ht}{H_{t}}
\newcommand{\Dt}{D_{t}}

\title{Recovering Time-Varying Correlation with DCC-GARCH:\\
\large A Controlled Study of Correlation Tracking and
Portfolio-Risk Misstatement}
\author{Eugen Soloviov}
\date{July 2026}

\begin{document}
\maketitle

\begin{abstract}
Sample correlation is an average over a window during which the true
dependence structure was moving. We ask, under known ground truth, how much a
dynamic model recovers that a static or rolling estimate throws away. We build
seeded synthetic data-generating processes in which the conditional
correlation follows a prescribed calm/crisis/recovery path (equicorrelation
$0.3$ in calm, ramping to $0.9$ in the crisis, partially recovering to $0.5$)
on top of GARCH$(1,1)$ volatilities, implement the two-step Dynamic Conditional
Correlation (DCC) estimator of Engle from scratch on top of the univariate
\texttt{arch} library, and measure recovery directly. First, DCC tracks the
correlation path with a crisis-window mean absolute error of $0.028$ on a
two-asset process, against $0.208$ for the static full-sample correlation ---
a factor of roughly $7$ --- while beating every fixed rolling window on
crisis-window error and on overall root-mean-square error; no single rolling
window is simultaneously best overall and in the crisis, and DCC is. On a
stationary DCC simulation the estimator recovers the true parameters
$(a,b)=(0.05,0.90)$ as $(0.049,0.914)$. Second, a one-step $95\%$ Value-at-Risk
backtest on an equal-weight portfolio shows the cost of ignoring correlation
dynamics: static covariance breaches at $0.111$ inside the crisis --- $2.2$
times the nominal $0.05$ --- while its full-sample breach rate of $0.049$ looks
perfectly calibrated, hiding the failure exactly when it matters; DCC holds the
crisis breach rate at $0.063$, close to the $0.060$ delivered by the true
covariance. Third, a DCC dynamic hedge ratio cuts hedged-spread variance by a
fraction $0.208$ inside the crisis relative to a static ordinary-least-squares hedge, and
tracks the true time-varying beta with a crisis error of $0.035$ against
$0.248$. Finally, a parameter-count comparison shows why DCC scales where full
multivariate GARCH cannot ($8$ versus $11$ versus $21$ free parameters at
$d=2$; $152$ versus $6275$ versus $3.25\times 10^{6}$ at $d=50$ for DCC, BEKK,
and VECH), and the estimated $(a,b)$ stay stable as the cross-section grows to
$d=10$. The deliverable is the calibrated method and the quantified cost of
static correlation, not a market claim: the processes are synthetic by design.
\end{abstract}

\section{Introduction}
\label{sec:intro}

Ask for the correlation between two assets and you typically get a single
number computed over some window nobody remembers choosing. That number
averages over a period in which the true dependence structure was moving
constantly. In calm markets assets drift apart; in a stress event they lock
together, and the diversification a portfolio was built on evaporates at the
moment it is most needed. Correlation is not a constant that occasionally gets
estimated badly --- it is a time series with its own clustering and regimes,
and treating it as a scalar is the multivariate cousin of assuming constant
volatility.

The econometrics literature offers a ladder of answers. Modeling the full
conditional covariance matrix directly --- VECH \citep{bollerslev1988vech} or
BEKK \citep{englekroner1995bekk} --- is general but its parameter count grows
quadratically or quartically in the number of assets $d$, confining it to a
handful of series in practice. The Constant Conditional Correlation model of
\citet{bollerslev1990ccc} factors the covariance into univariate volatilities
and a correlation matrix, buying tractability at the price of freezing the
correlation. Engle's Dynamic Conditional Correlation (DCC) model
\citep{engle2002dcc} keeps that factorization but lets the correlation matrix
evolve through a two-parameter recursion, estimated in two steps
\citep{englesheppard2001dcc}; \citet{cappiello2006asymmetric} add downside
asymmetry. DCC is attractive precisely because it scales: the volatility side
is $d$ independent univariate fits and the correlation side has only two free
parameters, whatever $d$ is.

What is easy to assert but harder to demonstrate is that this machinery
actually recovers a moving correlation better than the naive alternatives, and
that the difference matters for a downstream decision. Demonstrating it
requires knowing the truth. That is the gap this paper fills. We do not claim a
new estimator and we make no claim about real markets; the contribution is
calibration evidence under controlled ground truth --- synthetic, seeded
data-generating processes whose conditional correlation path we prescribe by
construction --- together with a self-contained DCC implementation built on the
univariate \texttt{arch} library \citep{sheppard_arch}, which fits univariate
GARCH \citep{engle1982arch,bollerslev1986garch} but has no multivariate layer.
Implementing the DCC layer is part of the contribution.

Concretely, we run four experiments (Section~\ref{sec:design}) against the
estimator of Section~\ref{sec:method}:
\begin{enumerate}
\item \textbf{Correlation tracking.} On two- and three-asset processes with a
  known calm/crisis/recovery correlation path, DCC tracks the path with lower
  crisis-window error than any fixed rolling window and far lower than static
  correlation, and on a stationary DCC simulation it recovers the true
  parameters $(a,b)$.
\item \textbf{Static-risk misstatement.} A one-step portfolio Value-at-Risk
  backtest shows that static and rolling covariance under-state risk precisely
  during the correlation spike, breaching far above the nominal rate inside the
  crisis, while DCC stays close to nominal.
\item \textbf{Dynamic hedge ratio.} A DCC time-varying hedge ratio reduces
  hedged-spread variance relative to a static hedge and tracks the true
  time-varying beta, most visibly inside the crisis.
\item \textbf{Dimensionality.} A parameter-count comparison quantifies DCC's
  linear scaling against BEKK's quadratic and VECH's quartic growth, and an
  empirical note confirms $(a,b)$ estimation stays stable as $d$ grows.
\end{enumerate}
This study accompanies a \href{https://marketmaker.cc}{marketmaker.cc} blog
post. All numbers derive from one seeded run of a public harness; a companion
script verifies every numeric claim in this manuscript against the run's JSON
output.

\section{The data-generating processes and ground truth}
\label{sec:dgp}

Every experiment draws from a synthetic process with a KNOWN conditional
covariance path, so that estimation error can be measured rather than argued
about. Let $r_{t}\in\mathbb{R}^{d}$ be the return vector at time $t$. We
generate
\begin{equation}
\label{eq:dgp}
r_{t} = \Dt\,L_{t}\,e_{t}, \qquad e_{t}\sim\mathcal{N}(0,I_{d}),
\qquad L_{t}L_{t}' = \Rt,
\end{equation}
where $\Dt=\diag(\sigma_{1,t},\dots,\sigma_{d,t})$ collects per-asset
conditional volatilities and $L_{t}$ is the Cholesky factor of the true
conditional correlation matrix $\Rt$. The true conditional covariance is
therefore $\Ht = \Dt\,\Rt\,\Dt$, known exactly at every step.

\paragraph{Correlation path.} $\Rt$ is an equicorrelation matrix with a single
time-varying off-diagonal $\rho_{t}$ following a piecewise path over
$T = 2000$ observations: a calm regime at $\rho = 0.3$, a linear ramp up to a
crisis level of $\rho = 0.9$, a crisis plateau, a recovery ramp down, and a
final calm regime at a partially recovered $\rho = 0.5$. The crisis window
--- the stress episode of elevated correlation --- spans observations $600$ to
$1300$. This is the ground-truth path every estimator in
Section~\ref{sec:res-track} is scored against.

\paragraph{Volatility path.} Each $\sigma_{i,t}^{2}$ follows a GARCH$(1,1)$
recursion $\sigma_{i,t}^{2} = \omega_{t} + \alpha\,\epsilon_{i,t-1}^{2} +
\beta\,\sigma_{i,t-1}^{2}$ with $\alpha = 0.10$ and $\beta = 0.85$, sharing a
variance intercept $\omega_{t}$ that rises from a calm $0.05$ to a crisis
$0.20$ over the same schedule as the correlation. Volatility and correlation
therefore spike together, as they do in real drawdowns: the crisis is a
high-volatility, high-correlation regime, and the resulting elevation of true
portfolio variance is what the risk experiment probes.

\paragraph{Stationary DCC process.} To test parameter recovery we also
generate standardized residuals directly from the DCC recursion itself
(Section~\ref{sec:method}) with fixed, known parameters $(a,b)$ and a fixed
target correlation, over $T = 4000$ observations. Here the ground truth is the
pair $(a,b)$, and the question is whether the estimator recovers it.

All random streams use \texttt{numpy}'s seeded generator, so every number is
bit-reproducible by one command.

\section{The DCC-GARCH estimator}
\label{sec:method}

\subsection{Factorization}

Following \citet{engle2002dcc}, the conditional covariance is factored as
\begin{equation}
\label{eq:factor}
\Ht = \Dt\,\Rt\,\Dt,
\end{equation}
with $\Dt$ the diagonal matrix of conditional standard deviations from $d$
univariate GARCH models and $\Rt$ the conditional correlation matrix. Elementwise,
$h_{ij,t} = \rho_{ij,t}\,\sigma_{i,t}\,\sigma_{j,t}$: the conditional covariance
of two assets is their dynamic correlation times each of their dynamic
volatilities. \citet{bollerslev1990ccc} froze $\Rt$ to a constant; DCC lets it
move.

\subsection{Step 1: univariate volatilities}

For each asset $i$ we fit a univariate GARCH$(1,1)$ with the \texttt{arch}
library \citep{sheppard_arch}, obtaining the conditional volatility
$\sigma_{i,t}$ and the standardized residual
\begin{equation}
\label{eq:zstd}
z_{i,t} = \frac{\epsilon_{i,t}}{\sigma_{i,t}}.
\end{equation}
Each $z_{i,t}$ has approximately unit conditional variance; whatever
co-movement remains in the vector $z_{t}=(z_{1,t},\dots,z_{d,t})'$ is
dependence, not a volatility artifact.

\subsection{Step 2: the correlation recursion}

We model an auxiliary process $Q_{t}$, a symmetric positive-definite matrix,
with a scalar GARCH-like recursion driven by the outer products of the
standardized residuals:
\begin{equation}
\label{eq:qrec}
Q_{t} = (1-a-b)\,\bar{Q} + a\,z_{t-1}z_{t-1}' + b\,Q_{t-1},
\end{equation}
where $\bar{Q}$ is the unconditional correlation matrix of the standardized
residuals --- plugged in as a sample estimate rather than optimized, the
\emph{correlation-targeting} device that confines the dimensionality to
$\bar{Q}$ and leaves only $a$ and $b$ to estimate. The constraints
$a,b>0$ and $a+b<1$ keep $Q_{t}$ positive definite and mean-reverting, exactly
analogous to a univariate GARCH$(1,1)$. Because $Q_{t}$ has a diagonal that is
not exactly one, we normalize it to a proper correlation matrix,
\begin{equation}
\label{eq:qnorm}
\Rt = \diag(Q_{t})^{-1/2}\,Q_{t}\,\diag(Q_{t})^{-1/2}, \qquad
\rho_{ij,t} = \frac{q_{ij,t}}{\sqrt{q_{ii,t}\,q_{jj,t}}},
\end{equation}
which is symmetric positive definite with unit diagonal at every step by
construction. The whole correlation layer has only two parameters, regardless
of whether $d=2$ or $d=50$.

\subsection{The quasi-log-likelihood}

Under $\epsilon_{t}\mid\mathcal{F}_{t-1}\sim\mathcal{N}(0,\Ht)$ the Gaussian
log-likelihood separates into a volatility part depending only on $\Dt$ and a
correlation part depending on $\Rt$ \citep{engle2002dcc}. Maximizing the
volatility part is exactly the $d$ univariate fits of Step~1; the correlation
part, given the standardized residuals, is the two-parameter objective
\begin{equation}
\label{eq:dccll}
\mathcal{L}^{C}(a,b) = -\frac{1}{2}\sum_{t=1}^{T}
\Big(\log|\Rt| + z_{t}'\,\Rt^{-1}\,z_{t}\Big),
\end{equation}
which we maximize over $(a,b)$ with a derivative-free multi-start search. It is
a \emph{quasi}-likelihood: the two-step estimator is consistent but not fully
efficient, and correct standard errors need the \citet{englesheppard2001dcc}
correction; for recovering the correlation path the point estimates are what
matter. Timing is strictly causal: in the recursion $\Rt$ is scored against
$z_{t}$ and only then is $Q$ advanced with $z_{t}z_{t}'$ for step $t+1$, so
$\Rt$ uses information through $t-1$ only --- the off-by-one that silently
inflates in-sample fit is avoided by construction.

\subsection{Baselines}

Against DCC we run the estimators a practitioner reaches for first. The
\emph{static} baseline is the full-sample Pearson correlation (or covariance),
a single matrix held fixed for all $t$. The \emph{rolling} baselines recompute
the Pearson correlation (or covariance) over a trailing window of the last $w$
observations, causal by construction, for $w\in\{30,60,120,250\}$. These are
exactly the estimators DCC must beat to justify its extra machinery.

\section{Experimental design}
\label{sec:design}

All experiments are generated and analyzed by one seeded Python harness
(\texttt{scripts/run\_all.py}, estimator in \texttt{scripts/dcc.py}) under
Python $3.14.6$ with NumPy $2.5.1$. A verification script
(\texttt{scripts/check\_paper\_numbers.py}) checks every numeric claim below
against \texttt{results/results.json} and fails on any mismatch; quantities are
rounded from that file.

\paragraph{Experiment 1: correlation tracking.} On the two- and three-asset
regime processes we compare the average pairwise correlation implied by DCC,
by each rolling window, and by the static estimate against the true path
$\rho_{t}$, reporting mean absolute error (MAE) and root-mean-square error
(RMSE) both overall and restricted to the crisis window. Separately, on the
stationary DCC process we fit $(a,b)$ and compare to the known truth.

\paragraph{Experiment 2: static-risk misstatement.} For an equal-weight
portfolio we compute one-step $95\%$ Gaussian Value-at-Risk,
$\mathrm{VaR}_{t} = z_{0.95}\sqrt{w'\Sigma_{t}w}$ with
$z_{0.95}\approx 1.645$, using in turn the static, rolling, DCC, and true
covariance for $\Sigma_{t}$, each built from information through $t-1$. A
breach is a realized portfolio loss exceeding $\mathrm{VaR}_{t}$. We report the
breach rate overall, in the crisis window, and in the calm periods; the nominal
rate is $0.05$.

\paragraph{Experiment 3: dynamic hedge ratio.} On the two-asset process we form
the minimum-variance hedge ratio of asset one against asset two,
$\beta_{t} = \rho_{t}\,\sigma_{1,t}/\sigma_{2,t}$, from DCC outputs, and compare
it to a static ordinary-least-squares hedge ratio (the full-sample covariance
ratio). Hedging with the lagged $\beta_{t-1}$ we measure the variance of the
hedged spread and its reduction relative to the static hedge, plus the tracking
error of the estimated $\beta_{t}$ against the true $\beta_{t}$.

\paragraph{Experiment 4: dimensionality.} We tabulate the free-parameter count
of full VECH, full BEKK, and DCC as $d$ grows, and re-run the parameter-recovery
experiment at $d\in\{2,3,5,10\}$ to check that the estimated $(a,b)$ stay stable
as the cross-section widens.

\section{Results}
\label{sec:results}

\subsection{Correlation tracking and parameter recovery}
\label{sec:res-track}

Table~\ref{tab:track} reports the tracking errors on the two-asset process.
The static full-sample correlation is a flat line at $0.673$, and it is
hopeless against a path that moves between $0.3$ and $0.9$: its crisis-window
MAE is $0.208$. DCC drives that down to $0.028$, a factor of roughly $7$. The
comparison against rolling windows is where the honest structure of the result
lives. No single fixed window is best everywhere: the window that tracks best
overall ($w=120$, overall MAE $0.060$) lags badly through the crisis (crisis
MAE $0.037$), while the window that tracks the crisis best ($w=60$, crisis MAE
$0.031$) is noisier overall. DCC is at or below the best window on every
slice at once --- crisis MAE $0.028$ (below the best window's $0.031$),
overall RMSE $0.083$ (below every window's), and overall MAE $0.061$
(indistinguishable from the best window's $0.060$). A short window is
responsive but noisy; a long window is smooth but stale; DCC's exponential,
adaptively weighted memory is both, which is exactly the tradeoff the two
parameters $(a,b)$ are there to resolve.

\begin{table}[t]
\centering
\caption{Experiment 1, two-asset process ($T = 2000$): tracking error of each
estimator against the true correlation path, overall and inside the crisis
window (observations $600$ to $1300$). DCC has the lowest crisis-window error
of all and the lowest overall RMSE; no single rolling window is best both
overall and in the crisis. The static estimate is the full-sample correlation
$0.673$.}
\label{tab:track}
\begin{tabular}{lrrrr}
\toprule
Estimator & MAE overall & RMSE overall & MAE crisis & RMSE crisis \\
\midrule
Static (full sample) & 0.238 & 0.259 & 0.208 & 0.217 \\
Rolling $w=30$ & 0.092 & 0.129 & 0.037 & 0.048 \\
Rolling $w=60$ & 0.069 & 0.100 & 0.031 & 0.040 \\
Rolling $w=120$ & 0.060 & 0.087 & 0.037 & 0.054 \\
Rolling $w=250$ & 0.071 & 0.102 & 0.074 & 0.101 \\
\textbf{DCC-GARCH} & 0.061 & 0.083 & \textbf{0.028} & \textbf{0.039} \\
\bottomrule
\end{tabular}
\end{table}

The three-asset process tells the same story: DCC's crisis-window MAE is
$0.066$, below every rolling window (best $0.070$ at $w=60$) and far below the
static $0.202$, and it fits $(a,b)=(0.031,0.966)$ --- the near-unit persistence
and small shock loading that is the canonical DCC fingerprint.

Parameter recovery is exact to within sampling noise. On the stationary
two-asset DCC process with true $(a,b) = (0.05,0.90)$, the estimator returns
$(0.049,0.914)$: an absolute error of $0.0015$ in $a$ and $0.014$ in $b$. On a
three-asset process with true $(0.04,0.92)$ it returns $(0.040,0.914)$. The
self-contained two-step estimator recovers the data-generating parameters, so
the tracking performance above is not an artifact of a mis-specified fit.

\subsection{Static covariance under-states portfolio risk in the spike}
\label{sec:res-risk}

Table~\ref{tab:risk} is the practical payoff. For the equal-weight two-asset
portfolio, the static full-sample covariance produces a $95\%$ VaR whose
overall breach rate is $0.049$ --- essentially perfect calibration, and
exactly the false comfort that makes static covariance dangerous. That average
hides a violent split: inside the crisis the static VaR is breached at $0.111$,
$2.2$ times the nominal $0.05$, while in the calm periods it breaches at only
$0.015$. A risk manager watching the full-sample number would see a
well-calibrated model while carrying more than double the intended tail risk
through every stress episode. The rolling covariance ($w=60$) improves the
crisis breach rate to $0.069$ but still runs hot, because a trailing window
that straddles the regime boundary is a blend of calm and crisis. DCC holds the
crisis breach rate to $0.063$, within touching distance of the $0.060$ that the
\emph{true} covariance itself delivers in finite sample --- the irreducible
floor. The three-asset portfolio repeats the pattern: static breaches at
$0.103$ in the crisis versus DCC's $0.070$ and the truth's $0.060$.

\begin{table}[t]
\centering
\caption{Experiment 2: one-step $95\%$ VaR breach rates for an equal-weight
portfolio (nominal $0.05$), by covariance estimator, overall and split into
crisis and calm periods. Static covariance looks calibrated overall ($0.049$)
but breaches at $0.111$ --- $2.2\times$ nominal --- inside the crisis; DCC
stays close to the true-covariance floor.}
\label{tab:risk}
\begin{tabular}{lrrr}
\toprule
Covariance estimator & Breach overall & Breach crisis & Breach calm \\
\midrule
Static (full sample) & 0.049 & 0.111 & 0.015 \\
Rolling ($w=60$) & 0.062 & 0.069 & 0.058 \\
DCC-GARCH & 0.056 & 0.063 & 0.052 \\
True covariance & 0.060 & 0.060 & 0.060 \\
\bottomrule
\end{tabular}
\end{table}

\subsection{A dynamic hedge ratio}
\label{sec:res-hedge}

On the two-asset process the static OLS hedge ratio is a constant $0.650$,
while the true minimum-variance hedge ratio averages $0.576$ and moves with the
regime. The DCC hedge ratio averages $0.575$ and tracks the true path with a
mean absolute error of $0.067$ overall and $0.035$ inside the crisis, against
the static hedge's $0.236$ and $0.248$. Trading the lagged hedge, the
DCC-hedged spread has a fractional variance reduction of $0.120$ relative to the static-hedged spread
overall, rising to $0.208$ inside the crisis --- precisely where the
static hedge, calibrated to the full-sample average, is most wrong. The dynamic
hedge does not create an edge; it keeps the traded spread market-neutral through
the regime shift instead of drifting into directional exposure exactly when
correlation moves.

\subsection{Dimensionality}
\label{sec:res-dim}

Table~\ref{tab:dim} is the structural reason DCC is used at all. Counting free
parameters, full VECH carries an intercept plus two coefficient matrices on the
half-vectorization of the covariance, growing as $O(d^{4})$; full BEKK carries a
triangular intercept plus two $d\times d$ coefficient matrices, $O(d^{2})$; DCC
carries three parameters per univariate GARCH plus the two correlation scalars,
$O(d)$. The gap is not asymptotic hand-waving: at $d=5$ it is $17$ parameters
for DCC versus $65$ for BEKK and $465$ for VECH, and at $d=50$ it is $152$
versus $6275$ versus $3.25\times 10^{6}$. VECH and BEKK are simply not
estimable on the cross-sections DCC handles routinely. The correlation layer
itself never grows past two parameters. Estimation stays well-behaved as the
cross-section widens: re-running the recovery of true $(a,b)=(0.05,0.90)$ at
$d=2,3,5,10$ returns $a$ within $0.006$ and $b$ within $0.015$ of the truth at
every width, with no drift as $d$ grows.

\begin{table}[t]
\centering
\caption{Experiment 4: free-parameter counts of full VECH, full BEKK, and DCC
(three per univariate GARCH$(1,1)$ plus the two correlation scalars) as the
number of assets $d$ grows. VECH is $O(d^{4})$, BEKK is $O(d^{2})$, DCC is
$O(d)$ with a correlation layer fixed at two parameters.}
\label{tab:dim}
\begin{tabular}{rrrr}
\toprule
$d$ & VECH & BEKK & DCC \\
\midrule
2 & 21 & 11 & 8 \\
3 & 78 & 24 & 11 \\
5 & 465 & 65 & 17 \\
10 & 6105 & 255 & 32 \\
20 & 88410 & 1010 & 62 \\
50 & 3252525 & 6275 & 152 \\
\bottomrule
\end{tabular}
\end{table}

\section{Discussion}
\label{sec:discussion}

\subsection{What the dynamic model buys}

The three substantive experiments are three views of one fact: a correlation
that moves must be tracked, not averaged. In tracking (Table~\ref{tab:track})
the static estimate's error is an order of magnitude worse than DCC's inside
the crisis, and the rolling windows force a choice --- responsiveness or
smoothness --- that DCC does not. In risk (Table~\ref{tab:risk}) the same
static average that looks calibrated in aggregate ($0.049$ against nominal
$0.05$) under-states the crisis tail by a factor of $2.2$, which is the failure
mode that matters: a risk model is not judged on its average behavior but on
its behavior in the tail, and the static model's tail arrives concentrated
inside exactly the window it cannot see. DCC's crisis breach rate of $0.063$
against the true-covariance floor of $0.060$ says the residual error is close
to irreducible. In hedging the dynamic ratio's crisis variance
reduction of $0.208$ is the same effect measured in a trading decision rather than a
diagnostic.

\subsection{Why correlation targeting and causal timing matter}

Two implementation choices carry the result and are worth isolating. Correlation
targeting --- plugging $\bar{Q}$ in as a sample estimate rather than optimizing
it --- is what collapses the correlation layer to two parameters and makes the
$d=50$ column of Table~\ref{tab:dim} possible; without it the correlation step
would itself be an $O(d^{2})$ optimization. Its cost, invisible in this study
by design, is that $\bar{Q}$ uses the full sample, so a strict walk-forward
evaluation must re-estimate it on training data only. Causal timing in the
recursion --- scoring $\Rt$ before advancing $Q$ --- is what keeps the tracking
numbers honest; the reverse convention would let $\Rt$ peek at the contemporaneous
shock and flatter every error in Table~\ref{tab:track}.

\subsection{The honest anatomy of the risk result}

The risk experiment is engineered to make one point cleanly, and its
limitations follow from that engineering. The crisis is a joint volatility and
correlation spike, so the static under-statement mixes two effects; we do not
decompose how much of the $2.2\times$ breach comes from stale volatility versus
stale correlation, only that ignoring the joint dynamics is what fails. The true
covariance breaches at $0.060$ rather than exactly $0.05$ because a
Gaussian-quantile VaR on a GARCH-mixture return is mildly miscalibrated in
finite sample even when the covariance is known --- which is the correct floor
to compare against, and DCC's $0.063$ essentially reaches it. These are single
seeded processes, not Monte Carlo frequency statements; they demonstrate the
mechanism and its magnitude under known truth, not its distribution across
markets.

\section{Limitations}
\label{sec:limitations}

\begin{itemize}
\item \textbf{Synthetic data, by design.} Both processes --- a prescribed
  correlation path over GARCH volatilities, and a stationary DCC recursion ---
  are chosen so that ground truth is known exactly, which is the point of a
  calibration study and also its boundary. Real returns have fat tails,
  asymmetric correlation responses, and structural breaks the Gaussian
  quasi-likelihood does not model; the experiments say what DCC recovers when
  its assumptions hold, not how it behaves when they do not.
\item \textbf{Scalar DCC dynamics.} One $a$ and one $b$ govern every pair.
  The asymmetric DCC of \citet{cappiello2006asymmetric}, which adds a
  downside-correlation term, is the natural extension for markets where
  crash-correlation dominates; we did not implement it.
\item \textbf{Correlation targeting leaks the full sample.} The in-sample
  $\bar{Q}$ used here is a look-ahead device acceptable in a controlled study
  but not in a backtest; the tracking and risk numbers are in-sample filtering
  results, not walk-forward performance.
\item \textbf{Gaussian quasi-likelihood.} We estimate under a Gaussian kernel;
  Student-$t$ innovations and the \citet{englesheppard2001dcc} standard-error
  correction would be needed for inference on real data. For recovering the
  point path they are not.
\item \textbf{Parameter counts are structural, not fitted.} Table~\ref{tab:dim}
  counts free parameters of the full VECH and BEKK specifications; the diagonal
  and scalar restrictions people actually use are smaller, but so is their
  generality, and none of them recover DCC's $O(d)$ scaling with a moving full
  correlation matrix.
\item \textbf{Single seeded runs.} The tracking, risk, and hedge experiments are
  worked examples with known truth, not frequency statements across seeds; the
  recovery experiment's stability across $d$ is the closest thing to a
  robustness sweep we report.
\end{itemize}

\section{Conclusion}
\label{sec:conclusion}

Under controlled conditions with a known moving correlation, DCC-GARCH does what
it promises. It tracks a calm/crisis/recovery correlation path with a
crisis-window error roughly a seventh of the static estimate's and at or below
the best fixed rolling window on every slice at once, and it recovers the
data-generating $(a,b)$ to within sampling noise. The cost of ignoring the
dynamics is concrete and one-sided: a static covariance that looks perfectly
calibrated in aggregate under-states portfolio tail risk by a factor of $2.2$
inside the correlation spike, while DCC tracks the true covariance to the
finite-sample floor; a static hedge drifts out of neutrality exactly when
correlation moves, where the dynamic hedge cuts spread variance by a fifth. And
because correlation targeting fixes the correlation layer at two parameters, all
of this scales to cross-sections where full multivariate GARCH cannot be
estimated at all. None of this is a claim about any market --- the processes are
synthetic and seeded --- but as a statement about the method, the dynamic model
recovers what the static one throws away, and the discarded information is most
valuable precisely when it is most costly to lack.

\paragraph{Reproducibility.} All code, tests, and outputs accompany this paper:
\texttt{scripts/run\_all.py} regenerates \texttt{results/results.json} from
fixed seeds (Python $3.14.6$, NumPy $2.5.1$); \texttt{scripts/dcc.py} contains
the self-contained DCC estimator; \texttt{scripts/check\_paper\_numbers.py}
verifies every numeric claim in this manuscript against that file and fails on
any mismatch; \texttt{tests/} contains deterministic invariant tests, including
that the estimator recovers a known $(a,b)$ within tolerance.

\bibliographystyle{plainnat}
\bibliography{refs}

\end{document}
